\documentclass[../../main.tex]{subfiles}
\begin{document}
\section{习题}\label{section:ch17-problems}

\begin{enumerate}
\item\label{exercise:ch17-01} 设 $M$ 是一个光滑流形。证明 $H^0_{\mathrm{dR}}(M)$ 的维数等于 $M$ 的连通分支数。
\item\label{exercise:ch17-02} 设 $(M,g)$ 是一个紧定向 Riemann 流形。利用习题~\cref{exercise:ch16-18}~中定义的 Hodge 星算子，定义形式上的 Laplace--Beltrami 算子 $\Delta\omega=\mathrm{d}\mathrm{d}^*\omega+\mathrm{d}^*\mathrm{d}\omega$。证明 $\Delta\omega=0$ 当且仅当 $\mathrm{d}\omega=0$ 且 $\mathrm{d}^*\omega=0$。
\item\label{exercise:ch17-03} 在上题基础上，证明：每个 de Rham 上同调类在紧 Riemann 流形上有唯一的调和代表元（Hodge 定理）。
\item\label{exercise:ch17-04} 计算 $H^p_{\mathrm{dR}}(\mathbb{R}^n\setminus\{0\})$。
\item\label{exercise:ch17-05} 计算 $H^p_{\mathrm{dR}}(\mathbb{T}^n)$。
\item\label{exercise:ch17-06} 计算 $H^p_{\mathrm{dR}}(\mathbb{R}\mathrm{P}^n)$。
\item\label{exercise:ch17-07} 设 $M$ 是一个连通紧光滑流形。证明 $H^n_{\mathrm{dR}}(M)\cong\mathbb{R}$ 当且仅当 $M$ 可定向。
\item\label{exercise:ch17-08} 设 $M$ 是一个紧连通光滑流形。证明任意光滑映射 $F:M\to S^n$ 的映射度是一个整数，并且与 $F$ 的光滑同伦类一一对应到同伦类集合。
\item\label{exercise:ch17-09} 证明对径映射 $A:S^n\to S^n$，$A(x)=-x$ 的映射度为 $(-1)^{n+1}$。
\item\label{exercise:ch17-10} 证明 $S^n$ 上不存在处处非零的连续切向量场，当 $n$ 为偶数时（毛球定理）。[提示：若存在，则可构造一个到对径映射的同伦，比较映射度。]
\item\label{exercise:ch17-11} 设 $M$ 是一个连通光滑流形。证明 $H^1_{\mathrm{dR}}(M)$ 同构于 $\operatorname{Hom}(\pi_1(M),\mathbb{R})$。
\item\label{exercise:ch17-12} 设 $U,V$ 是流形 $M$ 的两个开子集。写出 Mayer--Vietoris 序列在 $H^0$ 处的具体形式，并解释它的几何意义。
\item\label{exercise:ch17-13} 用 Mayer--Vietoris 序列计算环面 $\mathbb{T}^2$ 的所有 de Rham 群。
\item\label{exercise:ch17-14} 设 $M$ 是一个可定向紧光滑流形。证明 Poincaré 对偶在 de Rham 群层面给出 $H^p_{\mathrm{dR}}(M)\cong H^{n-p}_{\mathrm{dR}}(M)^*$。
\item\label{exercise:ch17-15} 设 $M$ 是一个光滑流形。证明 $M$ 的 de Rham 群只依赖于 $M$ 的同伦型。
\end{enumerate}
\end{document}
